\item El casco de un bote experimental se levantará sobre el agua mediante un aerodeslizador montado bajo su quilla, como se muestra en la figura P14.48. El aerodeslizador tiene forma de ala de avión. Su área proyectada sobre una superficie horizontal es $A$. Cuando el bote se remolca con una rapidez lo suficientemente alta, el agua se mueve en flujo laminar de modo que su rapidez de densidad promedio en la parte superior del aerodeslizador es $n$ veces mayor que su rapidez $v_b$ bajo el aerodeslizador. (a) Si ignora la fuerza de flotación, demuestre que la fuerza de sustentación hacia arriba que el agua ejerce sobre el aerodeslizador tiene una magnitud
$$ F \approx \frac{1}{2} (n^2 - 1) \rho v_b^2 A $$
(b) El bote tiene masa $M$. Demuestre que la rapidez de despegue es
$$ v \approx \sqrt{\frac{2Mg}{(n^2 - 1) \rho A}} $$
